%% =============================================================================
%% SEC 4 - SLOW-MOVING CAPITAL. ~980 words + Table 3. The centerpiece.
%% SOURCE: chapters/paper1_sec6_interdependencies_v2.tex, compressed ~6:1.
%% v2 additions: DCC monthly means (-0.29 / -0.10), GIRF decay (2-3 months, no
%% reverse response), VAR lag robustness -- all verified against the paper.
%% =============================================================================

Why does an alpha of 2--5\% per year survive in liquid, investment-grade
European markets? The slow-moving capital theory
\citep{duffie2010presidential, brunnermeier2009market, mitchell2012arbitrage}
gives a falsifiable answer: mispricings persist because the intermediary
capital that closes them is scarce, shared, and slow to redeploy. Our
three-strategy setting embeds an unusually sharp test. All three trades are
run by the same type of leveraged arbitrageur, but the CDS--bond and iTraxx
legs are \emph{held} by institutions, while the BTP~Italia bond is held by
unlevered retail households who face no margin calls and do not sell when
dealer capital withdraws \citep{vayanos2021preferred}. If a common capital
pool is the mechanism, the two institutional mispricings should co-move and
tighten their link in stress, while the retail-held basis should sit outside
the pool---and if anything move the other way, as forced institutional selling
in the two credit trades coincides with a flight of household demand toward
the linker. A common macroeconomic shock cannot generate that sign structure:
it would move all three claims alike, whoever holds them. We measure each
strategy's mispricing as the cross-sectional mean absolute basis over its
instrument universe,
\begin{equation}
m_{i,t}=\frac{1}{\lvert K_{i,t}\rvert}\sum_{k\in K_{i,t}}
\bigl\lvert b_{i,k,t}\bigr\rvert,
\end{equation}
where $K_{i,t}$ contains bonds with at least six months to maturity for the
two cash strategies (the negative bases only for the CDS--bond) and the four
on-the-run sub-indices for the skew, so that $\Delta m_{i,t}>0$ is a widening
of strategy $i$'s mispricing.
% SOURCE: sec6 intro + Section G identification logic.

Table~\ref{tab:smc}, Panel~A, shows the sign structure in the co-movement of
monthly mispricing changes. The institutional pair correlates at $+0.19$
unconditionally, and the correlation rises monotonically with stress, from
$-0.21$ in the LOW regime to $+0.32$ in HIGH---an increase significant at the
1\% level (Fisher $z$), and robust at the 10\% level to the
\citet{forbes2002no} correction for stress-driven volatility, so it is a
genuine tightening of the cross-market link rather than a variance artifact.
The BTP~Italia pairs behave as the experiment predicts: the correlation with
the CDS--bond basis is $-0.30$ on the full sample and deepens to $-0.47$ in
high stress. An interaction regression makes the asymmetry parametric:
\begin{equation}
\Delta m_{i,t} = \alpha + \beta_0\,\Delta m_{j,t}
              + \beta_1\!\left(\Delta m_{j,t}\times z_t\right)
              + \gamma\, z_t + \varepsilon_t,
\end{equation}
where $z_t$ is the contemporaneous standardized iTraxx Main spread; the level
term $\gamma z_t$ absorbs the widening common to both series, so that
$\beta_1$ isolates the incremental transmission under stress. Stress
\emph{amplifies} transmission within the institutional pair
($\beta_1=+0.08$, 5\% level) and \emph{deepens the negative} link toward the
retail-held basis ($\beta_1=-0.64$, 1\% level). Dynamic conditional
correlations tell the same story month by month: the BTP--CDS pair is negative
in every single month of the sample, averaging $-0.29$ against the CDS--bond
basis and $-0.10$ against the iTraxx skew---the mirror image of the
institutional pattern (Appendix~\ref{app:smc}).
% SOURCE: sec6 P1-P2; numbers Tables XIV-XV, Figure A.2 / DCC discussion.

Timing and persistence complete the mechanism. Granger tests on a VAR of the
three mispricings recover the liquidity pecking order of
\citet{mitchell2012arbitrage}: shocks originate in the most liquid,
standardized market---the iTraxx index, where dealer inventory adjusts
fastest---and propagate to the cash CDS--bond basis ($p=0.011$) and the
BTP~Italia basis ($p=0.007$) within a month, never the reverse; both links
survive at two and three lags. Generalized impulse responses trace the same
liquid-to-illiquid path with the same signs---positive into the CDS--bond
basis, negative into BTP~Italia---largest on impact and decaying to zero
within two to three months, with no significant response in the reverse
direction (Appendix~\ref{app:smc}). Persistence lines up with funding
dependence, the one margin on which the three trades differ mechanically: the
CDS--bond basis is the only trade that must finance a cash bond on a dealer
balance sheet for the life of the position, and its mispricing half-life
($\mathrm{HL}_i=\ln 2/\lvert\ln\varphi_i\rvert$ from an AR(1) fit to
$m_{i,t}$), 3.0
months, is one-third longer than the 2.2 months of the BTP~Italia and the
iTraxx skew---which need no such financing and are, on funding grounds,
indistinguishable in persistence.
% SOURCE: sec6 P3-P4; numbers Tables XVII-XVIII.

Panel~B closes the loop by estimating, strategy by strategy,
$\Delta m_{i,t} = a_i + \boldsymbol{\lambda}_i^{\prime}\mathbf{X}_t +
u_{i,t}$, where $\mathbf{X}_t$ collects observable proxies for the shadow
cost of intermediary funding. The Libor--OIS spread is
significant in all three equations and \emph{reverses sign across holder
bases}: $+16.63$ on the CDS--bond basis and $+3.85$ on the iTraxx skew, both
at the 1\% level, against $-15.99$ on BTP~Italia; market illiquidity repeats
the reversal. Dealer credit risk loads only on the one trade that must be
financed outright. A margin constraint binds only where the holder is
levered---an aggregate risk premium would move two claims on the same cash
flows in the same direction whoever holds them---so the sign reversal is the
signature of the funding channel itself. A placebo makes the falsification
explicit: the common factor of the three mispricings loads on the intermediary
proxies with the predicted signs (adjusted $R^2$ of 23.3\%), while on a
macroeconomic set---macro and real-activity uncertainty, long-term inflation
expectations, policy uncertainty---every coefficient is insignificant and the
adjusted $R^2$ falls by more than half, to 10.6\% (Appendix~\ref{app:smc}).
% SOURCE: sec6 P5 + Section G; numbers Table XIX.

Why, then, is the premium not competed away? The limits differ by trade but
share one root: the post-crisis cost of intermediary balance sheet. The
CDS--bond trade consumes repo capacity in a thin corporate-repo market subject
to recall risk \citep{bai2019cds, mitchell2012arbitrage}; the iTraxx skew
consumes regulatory capital---\citet{boyarchenko2016trends} document that as
required leverage ratios rose from 2--3\% to 5--6\%, the return on equity of
index-versus-constituent trades fell by a factor of two to three, producing
``new normal'' basis levels; and the BTP~Italia basis rests on a
retail holder base that neither monitors nor arbitrages it. The alpha of
Section~\ref{sec:alpha} is not an anomaly awaiting correction; it is the
equilibrium rent on a permanently scarcer balance sheet.
% SOURCE: sec6 Section G "Why the Mispricings Persist", compressed.

%% TABLE 3 - SLOW-MOVING CAPITAL AND THE HOLDER-BASE SIGN FLIP (compact).
%% Panel A = paper Table XIV; Panel B = paper Table XIX. All cells verified
%% against the compiled paper of 19 Jul 2026.
\begin{table}[t]
\centering
\caption{The holder-base experiment. Panel A: Pearson correlations of monthly
mispricing changes $\Delta m$, unconditional and by stress regime (iTraxx Main
5Y: LOW $<60$, MEDIUM $60$--$100$, HIGH $\geq 100$ bps); the LOW-to-HIGH rise
for the institutional pair is significant at 1\% (Fisher $z$, $p=0.005$) and
at 10\% after the Forbes--Rigobon correction ($p=0.070$). Panel B: each
$\Delta m$ regressed on proxies for the shadow cost of intermediary funding,
on the 148 common months; Newey--West $t$-statistics. Columns place the two institutionally
held trades first and the retail-held BTP~Italia last, so the sign reversal
reads across each row; Table~\ref{tab:rq3_pc1_funding} in the Appendix
reports the same regressions with full variable definitions. $^{***}$,
$^{**}$, $^{*}$: 1\%, 5\%, 10\%.
\label{tab:smc}}
\small
\begin{tabular}{lcccc}
\toprule
\multicolumn{5}{l}{\emph{Panel A: Co-movement of mispricing changes}}\\
Pair & Uncond. & LOW & MEDIUM & HIGH \\
\midrule
CDS--Bond vs.\ iTraxx (institutional pair) & $+0.19^{*}$ & $-0.21^{*}$ & $+0.20^{*}$ & $+0.32^{**}$ \\
CDS--Bond vs.\ BTP Italia (across bases)   & $-0.30^{***}$ & $-0.13$ & $-0.32^{***}$ & $-0.47^{**}$ \\
BTP Italia vs.\ iTraxx (across bases)      & $-0.10$ & $+0.15$ & $-0.15$ & $-0.14$ \\
\midrule
\multicolumn{5}{l}{\emph{Panel B: Response to intermediary funding conditions}}\\
Driver & \multicolumn{2}{c}{CDS--Bond} & iTraxx Skew & BTP Italia \\
\midrule
Libor--OIS spread        & \multicolumn{2}{c}{$+16.63^{***}$} & $+3.85^{***}$ & $-15.99^{**}$ \\
Amihud illiquidity       & \multicolumn{2}{c}{$+6.93^{**}$}   & $-0.42$       & $-20.51^{***}$ \\
Dealer CDS spread ($\Delta$5Y) & \multicolumn{2}{c}{$+0.11^{*}$} & $+0.02$   & $-0.06$ \\
HKM capital ratio        & \multicolumn{2}{c}{$+6.50$}        & $-0.67$       & $+8.81$ \\
Adj.\ $R^2$              & \multicolumn{2}{c}{0.165}          & 0.058         & 0.126 \\
\bottomrule
\end{tabular}
\end{table}

